\clearpage
\VerdierPhysicalPageAnchor{0242}%
\begingroup\small\setlength{\parskip}{0pt}\setlength{\emergencystretch}{1.5em}
\setlength{\abovedisplayskip}{4pt plus 2pt minus 2pt}%
\setlength{\belowdisplayskip}{4pt plus 2pt minus 2pt}%

\VerdierSemanticAnchor{ch4d-4-5-3-b-cont}
\begin{enumerate}
\item[1)] \textit{$X$ 是 $D^b(\mathcal A)$ 的对象。}
\item[2)] \textit{$X$ 是 $D^+(\mathcal A)$ 的对象，且 $F$ 右定常。}
\item[3)] \textit{$X$ 是 $D^-(\mathcal A)$ 的对象，且 $F$ 左定常。}
\item[4)] \textit{函子 $F$ 定常。}
\end{enumerate}

\VerdierSemanticAnchor{ch4d-4-5-3-proof}
这个命题立即由 (4.5.2) 与定义（第二章，4.4.2）得到。

\VerdierSemanticAnchor{ch4d-4-5-4}
\medskip\noindent\textbf{定理 4.5.4.}\quad
\textit{设
$F_1,F_2\colon D^*(\mathcal A)\longrightarrow\mathcal B$
（$*=+,-,b,{}$ ``空''）是两个上同调函子（相应地，
$G_1,G_2\colon D^*(\mathcal A)\longrightarrow D^{*'}(\mathcal B)$ 是两个
正合函子），并且 $m\colon F_1\longrightarrow F_2$（相应地，
$m\colon G_1\longrightarrow G_2$）是函子的态射（相应地，是正合函子的态射
（第一章，1.6.6.4））。由 $D^*(\mathcal A)$ 中所有满足如下条件的对象
$X$ 定义的 $D^*(\mathcal A)$ 的满子范畴：对每个整数 $p$，$m(X[p])$
都是同构；它是
$D^*(\mathcal A)$ 的严格满且饱和的三角子范畴（第二章，2.1.6），记作
$\operatorname{Is}(m)$。}

\VerdierSemanticAnchor{ch4d-4-5-4-a}
\noindent\textit{a) 若 $\mathcal A$ 的每个对象都是
$\operatorname{Is}(m)$ 的对象，则
$D^b(\mathcal A)\subset\operatorname{Is}(m)$。}

\VerdierSemanticAnchor{ch4d-4-5-4-b}
\noindent\textit{b) 若 $\mathcal A$ 的每个对象都是
$\operatorname{Is}(m)$ 的对象，且函子 $F_1$ 与 $F_2$（相应地，$G_1$ 与
$G_2$）右定常，则
$D^+(\mathcal A)\cap D^*(\mathcal A)\subset\operatorname{Is}(m)$。}

\VerdierSemanticAnchor{ch4d-4-5-4-c}
\noindent\textit{c) 若 $\mathcal A$ 的每个对象都是
$\operatorname{Is}(m)$ 的对象，且函子 $F_1$ 与 $F_2$（相应地，$G_1$ 与
$G_2$）左定常，则
$D^-(\mathcal A)\cap D^*(\mathcal A)\subset\operatorname{Is}(m)$。}

\VerdierSemanticAnchor{ch4d-4-5-4-d}
\noindent\textit{d) 若 $\mathcal A$ 的每个对象都是
$\operatorname{Is}(m)$ 的对象，且函子 $F_1$ 与 $F_2$（相应地，$G_1$ 与
$G_2$）定常，则
$D^*(\mathcal A)=\operatorname{Is}(m)$。}

\VerdierSemanticAnchor{ch4d-4-5-4-e}
\noindent\textit{e) 若存在 $\operatorname{Ob}(\mathcal A)$ 的子集 $M$，
使得：}

\noindent\textit{$\alpha$) $\mathcal A$ 的每个对象都同构于 $M$ 中某个
对象的商对象；}

\noindent\textit{$\beta$) $M$ 中的每个对象都是
$\operatorname{Is}(m)$ 的对象；}

\noindent\textit{并且函子 $F_1$ 与 $F_2$（相应地，$G_1$ 与 $G_2$）
左定常，则
$D^-(\mathcal A)\cap D^*(\mathcal A)\subset\operatorname{Is}(m)$。}

\VerdierSemanticAnchor{ch4d-4-5-4-f}
\noindent\textit{f) 若存在 $\operatorname{Ob}(\mathcal A)$ 的子集 $M$，
使得：}

\noindent\textit{$\alpha$) $\mathcal A$ 的每个对象都同构于 $M$ 中某个
对象的子对象；}

\noindent\textit{$\beta$) $M$ 中的每个对象都是
$\operatorname{Is}(m)$ 的对象；}

\noindent\textit{并且函子 $F_1$ 与 $F_2$（相应地，$G_1$ 与 $G_2$）
右定常，则
$D^+(\mathcal A)\cap D^*(\mathcal A)\subset\operatorname{Is}(m)$。}

\VerdierSemanticAnchor{ch4d-4-5-4-proof-start}
先注意，定理 (4.5.4) 中相应的括号内断言可由不带括号的断言推出：只须将函子
$G_1$ 与 $G_2$ 同典范上同调函子
$D^{*'}(\mathcal B)\longrightarrow\mathcal B$ 复合。证明 a)。设 $X$ 是
\endgroup
